\subsection{疲劳分析}
\label{sec:fatigue-var}

\subsubsection{取值}

可变径套管刮管工具的承力件分属两档材料，须建立两条 S-N 曲线。刮刀芯轴与心轴连接筒经设计方确认改用 4330V 的 1\,068/1\,180 档：$k_a=4.51\times1\,180^{-0.265}=0.692$、$k_b=0.748$、$k_c=0.85$、$k_d=1.0$、$k_e=0.814$，五项系数之积为 0.358，$S_e=0.5\times1\,180\times0.358\approx211.3$\,MPa；$1\times10^{3}$ 次处 $0.9\sigma_u=1\,062$\,MPa 尚低于屈服强度 1\,068\,MPa，无需截断，五点拟合得 $\sigma_a=5\,337\,N^{-0.2337}$，$\sigma_e=1\,062$\,MPa。上下接头仍为 4145H 的 760/965 档：$S_e=182.3$\,MPa，$0.9\sigma_u=868.5$\,MPa 超过屈服强度，按屈服截断取 760\,MPa，拟合得 $\sigma_a=3\,169\,N^{-0.2067}$。两条曲线的斜率相差 13\%，不得互换使用。

轴向拉伸工况的峰值 2\,156.6\,MPa 位于刮刀芯轴内孔变径台阶根部，属几何尖角引起的数值奇异。距该点 3$\sim$6\,mm 的探针读数 658.88\,MPa（节点 28315）与 660.2\,MPa（节点 28313）相差 0.2\%，沿轴向再远数毫米降至 277.9\,MPa（节点 28330），梯度之陡由此可见，取 658.88\,MPa 评定。该值可与解析式互证：按 $\phi58.05/\phi38.10$ 净截面（1\,506.5\,mm$^2$）反算名义应力为 663.8\,MPa，与有限元相差 0.7\%，说明模型在同一截面上给出的是可用的工程应力，而非数值假象。评定结果列于表~\ref{tab:fatigue_var}。

\begin{table}[htbp]
\centering
\caption{可变径套管刮管工具的疲劳评定结果（200 趟）}
\label{tab:fatigue_var}
\begin{tabularx}{\textwidth}{lCCC}
\hline
项目 & 轴向拉伸（$R=0$） & 抗扭（$R=-1$） & 判据 \\
\hline
额定工况 & 1\,000\,kN & 17\,kN$\cdot$m & --- \\
评定应力 / MPa & 658.88 & 921 & --- \\
应力幅与平均应力 / MPa & 329.44 & 921 / 0 & --- \\
等效对称应力幅 / MPa & 457.0 & 921 & --- \\
计算寿命 / 次 & $3.7\times10^{4}$ & $1.9\times10^{3}$ & --- \\
200 次累积损伤 & $5.4\times10^{-3}$ & 0.106 & $\le 1.0$ \\
疲劳安全系数 & 1.70 & 1.15 & $\ge 1.5$ \\
静强度安全系数 & 1.62 & 1.16 & $\ge 1.5$ \\
判定 & 满足 & 不满足 & --- \\
\hline
\end{tabularx}
\end{table}

拉伸工况的计算过程为：$\sigma_a=\sigma_m=329.44$\,MPa，$\sigma_{a,\mathrm{eq}}=329.44/(1-329.44/1\,180)=457.0$\,MPa，反解得 $N=(5\,337/457.0)^{1/0.2337}=3.7\times10^{4}$ 次，$D=200/3.7\times10^{4}=5.4\times10^{-3}$，$n_f=1/(329.44/1\,062+329.44/1\,180)=1.70$，$n_s=1\,068/658.88=1.62$，各项均在判据之内，与工具族其余三种 1\,000\,kN 级工具处于同一水平。

扭转工况的结论则不同。17\,kN$\cdot$m 下峰值 1\,832\,MPa 出现在刮刀芯轴 $\phi58$ 退刀槽根与键合面交界的凹角，属数值奇异；评定应力取退刀槽净截面上的 921\,MPa。该值与解析式相差 2\%：按薄壁圆轴反算得 $\tau=543.4$\,MPa、Mises 应力 941\,MPa，表明控制截面为退刀槽而非螺纹牙。按 $R=-1$，$\sigma_a=921$\,MPa、$\sigma_m=0$，无需 Goodman 折算，反解得 $N=(5\,337/921)^{1/0.2337}=1.9\times10^{3}$ 次，$D=200/1.9\times10^{3}=0.106$，$n_s=1\,068/921=1.16$，$n_f=1\,062/921=1.15$。静强度、疲劳安全系数与累积损伤三项同时不满足 1.5 与 1.0 的判据，为工具族内唯一的否决项。

其原因在截面而非材料：退刀槽槽底 $\phi58.05$ 较该端 Tr65$\times$4 螺纹的小径 $\phi60.5$ 尚小 2.45\,mm，净截面仅 1\,506.5\,mm$^2$，比螺纹根部（1\,734.6\,mm$^2$）小 13\%。该槽为螺纹退刀所必需，单独放大槽径将失去退刀作用，须连同该端螺纹一并放大，属设计变更。可考虑的处理有三条：其一，将该处净截面放大至 $\phi63.5$ 及以上（该端螺纹同步放大为 Tr68$\times$4），Mises 应力可降至 673\,MPa，安全系数回升至 1.59；其二，将额定扭矩降至 12.9\,kN$\cdot$m 以下，即该截面在 1.5 判据下允许的上限；其三，确认扭矩由中间筒体与连杆座分担的比例，按整机口径重新定值。三条之中以第一条改动最小，建议优先采用。

此外需说明两点。其一，本章模型只取轴向承力链上的五个零件，外筒、连杆座等传扭件全部抑制，扭矩全部由芯轴链承担，所得结果偏于保守；若整机试验或整机模型证明外筒分担了可观扭矩，该工况的安全系数将高于本节数值。其二，刮刀芯轴与心轴连接筒的零件图材料栏目前仍标注 4145H（$\ge$760\,MPa），本章按确认后的 1\,068\,MPa 计算，图纸须同步修改；若按图纸现值复核，螺纹根部安全系数仅为 1.14。扭转工况的应力云图与评定位置见第~\ref{sec:fea-var}~节与第~\ref{sec:torsion-var}~节，数值列于表~\ref{tab:torsion_var}。

轴向工况的控制部位为刮刀芯轴内孔变径台阶与螺纹根部，建议将过渡圆角加大至 R3 以上、成品螺纹根部进行磁粉探伤，并按 4330V/1\,068\,MPa 同步修改 BJG700004A、BJG700014A 两张零件图。
